\section*{Chapitre \ref{ch:b1:expression-control} --- Le contrôle de l'expression génétique}

\begin{solution}{exo:b1:expression-control:1}
\emph{lacI} (\omterm{def:b1:genomes:gene}{gène} du \omterm{def:b1:expression-control:operon}{répresseur}), promoteur, \omterm{def:b1:expression-control:operon}{opérateur}, \emph{lacZ},
\emph{lacY}, \emph{lacA}, et le site CAP. Sans lactose~: éteint
(\omterm{def:b1:expression-control:operon}{répresseur} fixé). Lactose et glucose présents~: à peine allumé
(\omterm{def:b1:expression-control:operon}{répresseur} détaché, mais pas de couple AMPc--CAP). Lactose sans
glucose~: pleinement allumé. Aucun des deux sucres~: éteint.
\end{solution}

\begin{solution}{exo:b1:expression-control:2}
Un produit qui agit en trans est une molécule diffusible qui peut agir
sur n'importe quelle copie de sa cible dans la \omterm{def:b1:cell-unit-of-life:cell}{cellule}~: le \omterm{def:b1:expression-control:operon}{répresseur}.
Un site qui agit en cis est une séquence d'\omterm{def:b1:nucleic-acids:chain}{ADN} qui n'affecte que les
\omterm{def:b1:genomes:gene}{gènes} qui lui sont physiquement liés~: l'\omterm{def:b1:expression-control:operon}{opérateur}.
\end{solution}

\begin{solution}{exo:b1:expression-control:3}
Latence d'environ \qty{36}{min} (de \qty{2}{h} à \qty{2.6}{h})~;
doublement en \qty{20}{min} sur glucose, en \qty{30}{min} sur lactose.
\end{solution}

\begin{solution}{exo:b1:expression-control:4}
\omterm{def:b1:expression-control:enhancer}{Amplificateur}~: séquence régulatrice, souvent lointaine, qui augmente la
transcription d'un \omterm{def:b1:genomes:gene}{gène} quand des facteurs s'y fixent. \omterm{def:b1:expression-control:enhancer}{Facteur de
transcription}~: \omterm{def:b1:proteins:peptide}{protéine} munie d'un domaine de liaison à l'\omterm{def:b1:nucleic-acids:chain}{ADN} et d'un
domaine d'activation ou de répression, qui règle la transcription. \omterm{def:b1:genomes:gene}{Gène}
de ménage~: \omterm{def:b1:genomes:gene}{gène} exprimé dans toute \omterm{def:b1:cell-unit-of-life:cell}{cellule} à un niveau constant.
\end{solution}

\begin{solution}{exo:b1:expression-control:5}
$I^-$~: constitutif. $O^c$~: constitutif. $I^s$~: non inductible.
$I^-/I^+$~: inductible (le \omterm{def:b1:expression-control:operon}{répresseur} $I^+$ agit en trans). $I^s/I^+$~:
non inductible ($I^s$ dominant). $I^+\,O^c\,Z^- / I^+\,O^+\,Z^+$~:
inductible --- le seul $Z$ en état de marche est derrière un \omterm{def:b1:expression-control:operon}{opérateur}
normal.
\end{solution}

\begin{solution}{exo:b1:expression-control:6}
Lac~: le sucre (l'inducteur) se lie au \omterm{def:b1:expression-control:operon}{répresseur} et le \emph{détache}
de l'\omterm{def:b1:expression-control:operon}{opérateur} --- les \omterm{def:b1:enzymes:enzyme}{enzymes} sont faites quand leur substrat est
présent. Trp~: l'\omterm{def:b1:proteins:aminoacid}{acide aminé} (le corépresseur) se lie au \omterm{def:b1:expression-control:operon}{répresseur} et
lui \emph{permet} de se lier --- les \omterm{def:b1:enzymes:enzyme}{enzymes} sont faites quand leur
produit est absent. Les voies cataboliques sont allumées par leur
entrée, les voies anaboliques éteintes par leur sortie.
\end{solution}

\begin{solution}{exo:b1:expression-control:7}
Sur glucose~: croissance normale (le glucose n'a pas besoin de CAP). Sur
lactose seul~: presque pas de croissance --- sans AMPc, la CAP ne peut
pas activer le promoteur lac, de sorte que l'\omterm{def:b1:enzymes:enzyme}{enzyme} reste à quelques
pour cent du niveau induit même avec le \omterm{def:b1:expression-control:operon}{répresseur} détaché. Sur les
deux~: croissance sur glucose, puis arrêt~; pas de seconde phase. \omterm{def:b1:enzymes:enzyme}{Enzyme}
faible dans tous les cas.
\end{solution}

\begin{solution}{exo:b1:expression-control:8}
L'atténuation exige que le \omterm{def:b1:gene-expression:ribosome}{ribosome} traduise le messager pendant que la
polymérase le transcrit encore, de sorte que la position du \omterm{def:b1:gene-expression:ribosome}{ribosome}
façonne l'\omterm{def:b1:nucleic-acids:chain}{ARN} situé derrière la polymérase. Chez les eucaryotes, la
transcription se fait dans le noyau et la traduction dans le \omterm{def:b1:eukaryotic-cell:organelle}{cytosol}, et
le message n'est traduit qu'une fois mûri et exporté.
\end{solution}

\begin{solution}{exo:b1:expression-control:9}
Les \omterm{def:b1:expression-control:enhancer}{amplificateurs} agissent à distance, dans l'une ou l'autre
orientation et de part et d'autre du \omterm{def:b1:genomes:gene}{gène}~; ils ne peuvent donc pas
agir en étant lus ni en plaçant directement la polymérase~; ils agissent
en fixant des facteurs qui touchent le promoteur par une boucle d'\omterm{def:b1:nucleic-acids:chain}{ADN},
ce pour quoi la distance et l'orientation importent peu.
\end{solution}

\begin{solution}{exo:b1:expression-control:10}
$5000\times 4096\times 4 = \num{8.2e7}$ \omterm{def:b1:water-small-molecules:atp}{ATP}~: \qty{0.8}{\%} du budget.
Taux de croissance abaissé de \qty{0.8}{\%}~: par génération, la part du
mutant est multipliée par $2^{-0.008} = 0.9945$. De $1:1$ à
$1:99$~: $0.9945^n = 1/99$, $n = \ln 99/0.00555 = 830$ générations ---
quelques semaines de culture continue.
\end{solution}

\begin{solution}{exo:b1:expression-control:11}
Chaque type cellulaire porte un jeu propre de \omterm{def:b1:expression-control:enhancer}{facteurs de transcription},
qui se lient aux \omterm{def:b1:expression-control:enhancer}{amplificateurs} de ses \omterm{def:b1:genomes:gene}{gènes} et recrutent des acétylases
qui en ouvrent la chromatine, tandis que les \omterm{def:b1:genomes:gene}{gènes} des autres types sont
méthylés et tassés dans une hétérochromatine qu'aucun facteur ne peut
atteindre. Les facteurs entretiennent mutuellement leur expression (un
réseau à états stables), et les marques de la chromatine sont copiées à
la réplication, de sorte que les filles héritent à la fois des facteurs
et des régions ouvertes et fermées. Le \omterm{def:b1:genomes:genome}{génome} est le même~; le \omterm{def:b1:genomes:genome}{génome}
accessible ne l'est pas, et l'inaccessibilité est héritée.
\end{solution}

\begin{solution}{exo:b1:expression-control:12}
Les contrôles d'une bactérie sont accordés au milieu~: un sucre, un
\omterm{def:b1:proteins:aminoacid}{acide aminé}, une température lie un \omterm{def:b1:expression-control:operon}{répresseur} ou un \omterm{prop:b1:expression-control:trp}{facteur sigma} et la
réponse est complète en quelques minutes, puis annulée dès que le signal
disparaît --- l'expression suit l'environnement. Une \omterm{def:b1:cell-unit-of-life:cell}{cellule} d'un corps
répond elle aussi à des signaux, mais ses décisions principales ont été
prises au cours du développement et verrouillées dans la chromatine~: ce
qu'elle peut exprimer a été fixé par les facteurs dont elle a hérité et
par les marques de son \omterm{def:b1:nucleic-acids:chain}{ADN}, changées seulement par la division et pour
l'essentiel irréversibles. La bactérie est une lectrice du présent~; la
\omterm{def:b1:cell-unit-of-life:cell}{cellule} différenciée a une mémoire, et cette mémoire est son identité.
\end{solution}

\begin{solution}{pb:b1:expression-control:1}
\textbf{1.} $0.5\times 10^{-3}\times 180 = \qty{0.09}{g/L} =
\qty{9e-5}{g/mL}$.
\textbf{2.} $9\times 10^{-5}/1.5\times 10^{-12} = \num{6e7}$ divisions.
\textbf{3.} $10^7 + \num{6e7} = \num{7e7}$ \omterm{def:b1:cell-unit-of-life:cell}{cellules} par mL.
\textbf{4.} $7 = 2^n$~: $n = 2.8$ doublements, soit \qty{56}{min}.
\textbf{5.} Éteint~: le lactose (l'allolactose) a détaché le \omterm{def:b1:expression-control:operon}{répresseur},
mais le glucose maintient l'AMPc bas, de sorte que la CAP n'est pas
fixée et que le promoteur est presque muet~; de plus la perméase est
rare, si bien que peu de lactose entre.
\textbf{6.} L'AMPc monte et se lie à la CAP, qui se lie au promoteur.
\textbf{7.} Le contrôle négatif était levé mais le contrôle positif
manquait~: un promoteur faible sans CAP démarre rarement (\omterm{def:b1:expression-control:cap}{répression
catabolique}).
\textbf{8.} $5000/1800 = 2.8$ tétramères par seconde~; $2.8\times 4096 =
\num{11400}$ résidus par seconde.
\textbf{9.} Capacité $20\,000\times 20 = \num{4e5}$ résidus par
seconde~: \qty{3}{\%} des \omterm{def:b1:gene-expression:ribosome}{ribosomes}.
\textbf{10.} $5000\times 4096\times 110\times 1.66\times 10^{-24} =
\qty{3.7e-15}{g}$~: \qty{2.5}{\%} des \omterm{def:b1:proteins:peptide}{protéines} de la \omterm{def:b1:cell-unit-of-life:cell}{cellule}.
\textbf{11.} $5000/5 = 1000$.
\textbf{12.} Le lactose ne peut pas traverser la membrane sans la
perméase~; un mutant \emph{lacY}$^-$ a l'\omterm{def:b1:enzymes:enzyme}{enzyme} mais pas de substrat à
l'intérieur et ne croît pas sur lactose (et s'induit mal, puisque
l'inducteur est fabriqué à partir du lactose dans la \omterm{def:b1:cell-unit-of-life:cell}{cellule}).
\textbf{13.} $1.0\times 10^{-3}\times 342 = \qty{0.342}{g/L}$, qui
valent $2\times 180 = \qty{360}{g}$ de glucose par mole, soit
\qty{0.36}{g/L} $= \qty{3.6e-4}{g/mL}$.
\textbf{14.} $3.6\times 10^{-4}/1.5\times 10^{-12} = \num{2.4e8}$
divisions~; total final $\num{7e7} + \num{2.4e8} = \num{3.1e8}$ par mL.
\textbf{15.} $3.1/0.7 = 4.4 = 2^n$~: $n = 2.1$ doublements, soit
\qty{64}{min}.
\textbf{16.} Glucose \qty{56}{min}, latence \qty{30}{min}, lactose
\qty{64}{min}~: \qty{150}{min} en tout.
\textbf{17.} Le lactose doit d'abord être hydrolysé, étape enzymatique
supplémentaire dont la vitesse limite l'approvisionnement en glucose~;
et la perméase emploie le gradient de protons pour l'importer, coût que
le transport du glucose ne paie pas.
\textbf{18.} $I^-$~: \omterm{def:b1:enzymes:enzyme}{enzyme} à \num{5000} tout du long (constitutif, une
fois le glucose épuisé~; sur glucose la CAP la limite encore, de sorte
que le niveau est intermédiaire)~; sa latence est plus courte ou nulle,
de sorte que la courbe montre une pause plus faible~; le coût le ralentit
légèrement sur glucose.
\textbf{19.} CAP$^-$~: croît sur glucose jusqu'à $\num{7e7}$, puis
s'arrête~: l'\omterm{def:b1:expression-control:operon}{opéron lactose} ne peut pas être activé, l'\omterm{def:b1:enzymes:enzyme}{enzyme} reste vers
5 molécules, et le lactose n'est jamais employé.
\textbf{20.} Sur lactose seul~: aucune différence de taux d'\omterm{def:b1:enzymes:enzyme}{enzyme} une
fois induit (les deux pleinement allumés) et pratiquement aucune en
croissance~; l'$O^c$ n'évite que le délai d'induction. Sur glucose
seul~: l'$O^c$ fabrique l'\omterm{def:b1:enzymes:enzyme}{enzyme} inutilement (dans la mesure où la CAP
le permet), le type sauvage non.
\textbf{21.} $5000\times 4096\times 4 = \num{8.2e7}$ \omterm{def:b1:water-small-molecules:atp}{ATP}~: \qty{0.8}{\%}
de $10^{10}$.
\textbf{22.} Temps de génération allongé de \qty{0.8}{\%}~: taux de
croissance abaissé de \qty{0.8}{\%}, $\delta = 0.008$ doublement par
génération.
\textbf{23.} $r = 2^{-0.008} = 0.99447$~; $r^n = 0.111$~: $n =
\ln 0.111/\ln 0.99447 = 2.20/0.00555 = 400$ générations.
\textbf{24.} Sur lactose, l'\omterm{def:b1:enzymes:enzyme}{enzyme} est de toute façon nécessaire, de
sorte que le mutant ne paie rien et gagne la latence~: il n'est pas
contre-sélectionné et il l'est parfois positivement~; et les souches de
laboratoire sont cultivées pour la commodité, non pour la compétition.
\textbf{25.} Facteur d'induction 1000 (de 5 à \num{5000} molécules)~; la
synthèse inutile coûte \qty{0.8}{\%} du budget \omterm{def:b1:water-small-molecules:atp}{ATP}~; l'expérience dure
environ deux heures et demie.
\end{solution}
